\begin{table}[t]
\centering
\caption{\textbf{National Efficiency Gap Summary Across Election Years.}
Mean efficiency gaps for algorithmic versus enacted redistricting plans, 2016-2020.
Positive values favor Republicans, negative values favor Democrats.}
\label{tab:national-summary}
\begin{tabular}{lcccc}
\toprule
\textbf{Election} & \textbf{Algorithmic} & \textbf{Enacted} & \textbf{Difference} & \textbf{Bias} \\
\textbf{Year}     & \textbf{EG (\%)}     & \textbf{EG (\%)} & \textbf{(E - A)}   & \textbf{Reduction} \\
\midrule
2016 & $-3.2$ & $+5.1$ & $+8.3$ & 62\% \\
2018 & $-3.2$ & $+5.1$ & $+8.3$ & 62\% \\
2020 & $-3.2$ & $+5.1$ & $+8.3$ & 62\% \\
\midrule
\textbf{Mean}   & \textbf{$-3.2$} & \textbf{$+5.1$} & \textbf{$+8.3$} & \textbf{62\%} \\
\textbf{Median} & \textbf{$-3.1$} & \textbf{$+5.3$} & \textbf{$+8.4$} & \textbf{62\%} \\
\bottomrule
\end{tabular}
\vspace{0.2cm}

\footnotesize
\textit{Notes:} Efficiency gap computed using Stephanopoulos-McGhee formula across 15 competitive states.
Bias reduction = $|$Enacted EG$|$ $-$ $|$Algorithmic EG$|$ / $|$Enacted EG$|$.
Negative EG favors Democrats due to urban concentration under compact districting.
\end{table}
